% Standalone research entry 213; batch 208--217.
% Original section material preserved from Pasted text(3).txt.
% Research additions: 27 September 2026.
% Compile twice with: lualatex 213.tex
% !TeX program = lualatex
% Compile twice: lualatex open_problems_500.tex
% All references and prose are embedded; no BibTeX database or external inputs.
\documentclass[11pt,a4paper]{article}
\usepackage[margin=24mm,headheight=14pt]{geometry}
\usepackage{fontspec}
\setmainfont{Latin Modern Roman}
\setsansfont{Latin Modern Sans}
\setmonofont{Latin Modern Mono}
\usepackage{amsmath,amssymb,mathtools}
\usepackage{unicode-math}
\setmathfont{Latin Modern Math}
\usepackage{microtype}
\usepackage{enumitem,xurl,needspace}
\usepackage[dvipsnames]{xcolor}
\usepackage{hyperref}
\hypersetup{unicode=true,colorlinks=true,linkcolor=MidnightBlue,urlcolor=MidnightBlue,
 pdftitle={500 Mathematical Problem Statements},pdfauthor={Source-annotated compilation},bookmarksnumbered=true}
\usepackage{bookmark}
\usepackage{tocloft}
\setlength{\cftsecnumwidth}{3em}
\cftsetpnumwidth{2.5em}
\cftsetrmarg{4em}
\usepackage{titlesec}
\titleformat{\section}{\Large\bfseries\raggedright}{\thesection}{0.7em}{}
\usepackage{fancyhdr}
\pagestyle{fancy}\fancyhf{}
\fancyhead[L]{\small\sffamily Mathematical problem statements}
\fancyhead[R]{\small Source edition 22}
\fancyfoot[C]{\thepage}
\setlength{\parindent}{0pt}
\setlength{\parskip}{5pt plus 1pt}
\setlength{\emergencystretch}{3em}
\setcounter{tocdepth}{1}
\setcounter{secnumdepth}{1}
\setlist[itemize]{leftmargin=3em,itemsep=5pt,topsep=4pt,parsep=0pt}
\allowdisplaybreaks
\newcommand{\N}{\mathbb{N}}
\newcommand{\Z}{\mathbb{Z}}
\newcommand{\Q}{\mathbb{Q}}
\newcommand{\R}{\mathbb{R}}
\newcommand{\C}{\mathbb{C}}
\newcommand{\F}{\mathbb{F}}
\newcommand{\A}{\mathbb{A}}
\newcommand{\PP}{\mathbb P}
\newcommand{\E}{\mathbb{E}}
\newcommand{\eps}{\varepsilon}
\newcommand{\dd}{\,\mathrm d}
\newcommand{\sref}[2]{\hyperlink{source:#1:#2}{[S#2]}}
\newcommand{\eref}[2]{\hyperlink{extra:#1:#2}{[E#2]}}
% Turn the next switch off for a shorter reading copy. The quotations preserve
% every exactTarget field, including qualifications omitted from a short summary.
\newif\ifshowcatalogue
\showcataloguetrue
\newcommand{\cataloguescope}[1]{\ifshowcatalogue
 \paragraph{Exact scope in the supplied catalogue.}
 \begin{quote}\small #1\end{quote}\fi}

% Standalone wrapper; no external inputs or bibliography files are required.
\fancyhead[L]{\small\sffamily Research notebook: section 213}
\fancyhead[R]{\small 27 September 2026}
\hypersetup{pdftitle={Research attempt 213},pdfauthor={Research notebook}}
\setlength{\parskip}{4pt plus 1pt}
\setlist[itemize]{before=\small\interlinepenalty10000,itemsep=3pt}
\begin{document}
\setcounter{section}{212}
% ================================================================
% problemId: problem.polynomial-diameter-bound-for-polytopes
% source releaseRank: 213; source releaseStatus: open
% exactTarget SHA-256: dc658c9c77171afa204e6501eaf44cdda8900867b9a2a6aa05721c72b50d4e34
\clearpage
\section{\texorpdfstring{Polynomial Diameter Bound for Polytopes}{Polynomial Diameter Bound for Polytopes}}\label{213}
\subsection{Definitions and mathematical statement}
A pointed polyhedron is an intersection of finitely many closed affine half-spaces in $\R^d$ containing no affine line; assume its affine dimension is $d$. Its graph has vertices the zero-dimensional faces and edges joining pairs that share a bounded one-dimensional face. Graph diameter is the maximum shortest edge-path distance between vertices. The target is a single polynomial $p$ such that
\[
 \operatorname{diam}(P)\le p(n,d)
\]
for every pointed $d$-dimensional polyhedron with $n$ facets, including unbounded polyhedra. Coordinates and their bit lengths do not enter the bound. This asks for a polynomial diameter estimate, not the stronger historical linear bound nor a polynomial-time pivot rule.
\par\smallskip\textit{Formulation and concept references:} \sref{213}{1}.
\cataloguescope{Prove or disprove that there is a polynomial p such that the vertex-edge graph of every pointed d-dimensional polyhedron, bounded or unbounded, with n facets has diameter at most p(n,d).}
\subsection{Short English statement}
Can the graph distance between any two vertices of a pointed polyhedron always be bounded by a polynomial in its dimension and number of facets?
% BEGIN RESEARCH ADDITION 213
\subsection{Research attempt}

\paragraph{Current frontier.}
Kalai--Kleitman give a quasipolynomial worst-case bound rather than a
polynomial with a fixed exponent \eref{213}{1}. The catalogue's formulation
reference studies random polytopes; its probabilistic conclusions do not
bound every pointed polyhedron \sref{213}{1}. Particularly relevant adjacent
progress is Natura's February 2026 preprint: its Theorem 1.1 bounds
\emph{circuit} diameter by $O(m^2\log m)$ in standard equality form with
$m$ independent equations \eref{213}{2}. Circuit walks may visit nonvertices,
and the result is not a bound for the graph in this entry. The attempt below
investigates the precise conversion needed, rather than identifying these
two notions of diameter.

\paragraph{Attack route.}
A direct attack would find a polynomial-length edge path. A potentially more
flexible route is to begin with a short continuous or circuit walk and replace
it by a chain of low-dimensional faces, inside which vertex enumeration has
polynomial cost. The missing geometric lemma would have to keep the face
dimension uniformly bounded, or provide a different polynomial navigation
bound inside the visited faces. I prove the conversion statement and then
try to derive its hypothesis from support-minimality of circuit directions.
The latter derivation fails in an explicit family.

\paragraph{Attempt.}
\textit{1. Standard form does not discard the unbounded case.}
Write a full-dimensional pointed polyhedron as
$P=\{x\in\R^d:Bx\le b\}$ with $n$ facet inequalities. Its normals span
$\R^d$: otherwise a vector in $\ker B$ would produce an affine line in $P$.
The slack map $x\mapsto b-Bx$ is therefore an affine injection and identifies
$P$ with an affine subspace intersected with $\R^n_{\ge0}$, equivalently
\[
 Q=\{s\in\R^n:As=c,\ s\ge0\},\qquad\operatorname{rank}A=n-d.
\]
Faces, vertices, bounded edges, and their incidences are preserved. No
boundedness or rationality assumption on the coefficients was used. In the
remainder, $m$ denotes the equation rank in standard form, not the original
dimension $d$.

\textit{2. Face-chain lifting lemma.}
Let $u,v$ be vertices of a pointed polyhedron with at most $n$ facets.
Suppose faces $F_1,\ldots,F_L$ satisfy
\[
 u\in F_1,\quad v\in F_L,\quad
 F_i\cap F_{i+1}\ne\varnothing,\quad \dim F_i\le r.
\]
Then
\[
 \operatorname{dist}_{\mathrm{graph}}(u,v)
 \le L\left(\sum_{j=0}^{r}\binom nj-1\right).                  \tag{213.1}
\]
For $r=0$ the hypotheses force $u=v$, so the zero right-hand side is correct.
To prove the other cases, every nonempty face of a pointed polyhedron has
a vertex. Choose a vertex $w_i$ of $F_i\cap F_{i+1}$, with $w_0=u$ and
$w_L=v$. A vertex of a face is a vertex of the entire polyhedron. A
$k$-dimensional face described in its affine hull by at most $n$ inequalities
has at most $\binom nk$ vertices: a vertex is specified by some $k$
independent active inequalities, and a fixed such choice determines at most
one point. Degeneracy only permits several choices for the same vertex and
does not invalidate the upper bound.

For completeness the vertex graph of each such face is connected even if
the face is unbounded. Choose a linear functional strictly positive on its
nonzero recession directions and generic on its finite vertex set. It has
a unique minimizing vertex. From any other vertex there is a descending
incident edge: the tangent cone is generated by edge directions and
recession rays, and descent cannot be supplied solely by recession rays
on which the functional is positive. Strict descent visits finitely many
vertices and terminates at the minimizer. Hence there is a simple graph
path between $w_{i-1}$ and $w_i$ inside $F_i$, of length at most the number
of its vertices minus one. Summing proves (213.1). All used edges of a face
are also edges of the original polyhedron.

A piecewise linear walk whose successive segments lie in faces of dimension
at most $r$ gives such a chain. Consequently, the following precise
hypothesis would imply the full target:
\begin{equation}\tag{213.2}
 \begin{gathered}
 \text{A fixed integer }r\text{ and a polynomial }L(n,d)\text{ exist such that}\\
 \text{every vertex pair admits a chain of at most }L(n,d)\\
 \text{intersecting faces, each of dimension at most }r.
 \end{gathered}
\end{equation}
This is a sufficient geometric reformulation, not a bound already obtained.
In fact polynomial graph diameter itself would give (213.2) with $r=1$;
the interest would be a new proof of a bounded-width chain by non-edge
methods. Allowing $r$ to grow like $\log n$ in (213.1) only recovers a
quasipolynomial scale, so the uniformity in $r$ cannot be hidden in notation.

\textit{3. A maximal circuit step need not end at a vertex.}
Take
\[
 A=\begin{pmatrix}1&0&1&1\\0&1&1&-1\end{pmatrix},\quad
 b=\binom11,\quad Q=\{x\ge0:Ax=b\}.
\]
The point $u=(1,1,0,0)$ is a vertex, and
$g=(-2,0,1,1)$ lies in $\ker A$. Its support $\{1,3,4\}$ is minimal:
every pair of those columns is linearly independent. The largest feasible
step is $\lambda=1/2$, giving
\[
 u+\tfrac12 g=(0,1,1/2,1/2)
 =\tfrac12(0,0,1,0)+\tfrac12(0,2,0,1).                        \tag{213.3}
\]
The two points on the right are distinct vertices. The endpoint on the left
is not a vertex. Thus the proposed conversion ``each maximal circuit step
is an edge step'' is false for a two-dimensional bounded polytope already.
The midpoints of the segment from $u$ to this endpoint have every coordinate
positive, so that segment runs through the relative interior of $Q$.

\textit{4. The obstruction has arbitrarily large face width.}
There is a direct extension of (213.3). Fix any $m\ge2$, let
\[
 w_1=e_1,\ldots,w_{m-1}=e_{m-1},\quad
 w_m=-\sum_{j=1}^{m-1}e_j\quad\text{in }\R^{m-1},
 \qquad c_j=(1,w_j)\in\R^m,
\]
and set
\[
 A_m=[\,I_m\mid c_1\ \cdots\ c_m\,],\quad
 Q_m=\{x\ge0:A_mx=\mathbf1\}.
\]
The columns $c_j$ are linearly independent. Indeed, a relation between them
has coefficient sum zero, whereas its remaining coordinates force all
coefficients to be equal. The relation
$\sum_jc_j=m e_1$ is the unique relation, up to scale, between the $m+1$
columns $e_1,c_1,\ldots,c_m$, and all of its coefficients are nonzero.
It therefore defines a circuit direction
\[
 g=(-m,0,\ldots,0;1,\ldots,1).
\]
Starting at the vertex $u=(\mathbf1;\mathbf0)$, the maximal step has size
$1/m$. At every intermediate time $0<t<1/m$ all $2m$ coordinates are
positive. Thus the segment's minimal containing face is \emph{all} of
$Q_m$, of dimension $2m-m=m$. At the endpoint only coordinate $1$ is zero;
the remaining columns still span $\R^m$, so its minimal face has dimension
$(2m-1)-m=m-1$. In particular it is not a vertex.

These polyhedra are pointed automatically, being contained in an orthant.
They are even bounded: the first equation bounds $x_1$ and the sum of the
last $m$ coordinates by one; the other equations then bound the remaining
basic coordinates. A positive intermediate point verifies the asserted
affine dimension. The phenomenon is thus not caused by recession rays,
degeneracy, or enormous coordinate bit lengths. A circuit's support bound
$m+1$ does not bound by an absolute constant the dimension of the face
traversed by that circuit. Exact rank checks for $m=2,\ldots,8$ accompany
the general proof, but the proof, not this finite list, establishes the
arbitrarily large dimension.

\paragraph{Stress test.}
The examples do not have large graph diameter and are not counterexamples
to the requested polynomial bound. They disprove a conversion lemma, not
the conclusion. Nor do they refute the circuit theorem in \eref{213}{2};
its steps are expressly allowed to have the behavior above. A more careful
route could replace one large-face segment by several low-face segments,
but no polynomial bound on the replacement cost has been supplied.

The lifting proof uses vertices of face intersections, not arbitrary nearest
vertices; nearest-point rounding would not preserve adjacency. It also
explicitly includes unbounded faces. Truncating an unbounded polyhedron and
then forgetting the new facets would not be a valid alternative without
controlling the introduced paths and combinatorial parameters. Finally,
short paths between two \emph{specified} vertices do not themselves give a
pivot rule that finds an unknown optimal vertex. Nothing here implies a
strongly polynomial linear-programming algorithm or makes an unsupported
complexity-theoretic inference.

\paragraph{Outcome.}
\textbf{Outcome: Conditional reduction.}

A bounded-face-width chain of polynomial length would yield a polynomial
graph-diameter bound by (213.1), with all unbounded and degenerate cases
included. The serious obstruction to obtaining such chains from short
circuit walks is demonstrated by a complete family of maximal circuit steps
through faces of unbounded dimension. The required refinement theorem, or
a different polynomial navigation estimate for these high-dimensional
faces, remains unproved. No unconditional improvement of the worst-case
graph-diameter bound is claimed.
% END RESEARCH ADDITION 213
\subsection{Sources}
\begin{itemize}\raggedright
\item[\textnormal{[S1]}]\hypertarget{source:213:1}{} Bonnet, Dadush, Grupel, Huiberts, and Livshyts, SoCG 2022.
\url{https://drops.dagstuhl.de/entities/document/10.4230/LIPIcs.SoCG.2022.18}.
{\small\textit{Catalogue formulation source.}}
% BEGIN RESEARCH REFERENCES 213
\item[\textnormal{[E1]}]\hypertarget{extra:213:1}{} Gil Kalai and Daniel J. Kleitman,
\emph{A quasi-polynomial bound for the diameter of graphs of polyhedra},
Bulletin of the American Mathematical Society 26 (1992), 315--316;
arXiv:math/9204233, main theorem.
\url{https://arxiv.org/abs/math/9204233}.
\item[\textnormal{[E2]}]\hypertarget{extra:213:2}{} Bento Natura,
\emph{Circuit Diameter of Polyhedra is Strongly Polynomial},
arXiv:2602.06958v2 (10 February 2026), Section 1.1, Theorem 1.1,
Theorem 3.1, and Corollary 3.2.
\url{https://arxiv.org/html/2602.06958v2}.
{\small\textit{Recent preprint concerning circuit, not vertex-edge, diameter; not independently proof-audited here.}}
% END RESEARCH REFERENCES 213
\end{itemize}

\end{document}
